\documentclass[12pt,a4paper]{article}

%% ---- Encodage et langue ----
\usepackage[utf8]{inputenc}
\usepackage[T1]{fontenc}
\usepackage[francais]{babel}

%% ---- Mathématiques ----
\usepackage{amsmath}
\usepackage{amssymb}
\usepackage{amsthm}

%% ---- Mise en page ----
\usepackage{geometry}
\geometry{left=2.5cm, right=2.5cm, top=2.5cm, bottom=2.5cm}
\usepackage{parskip}
\setlength{\parindent}{0pt}

%% ---- En-tête / pied de page ----
\usepackage{fancyhdr}
\pagestyle{fancy}
\fancyhf{}
\rhead{Analyse de Régression}
\lhead{Cerveau KATCHON}
\cfoot{\thepage}

%% ---- Commandes personnalisées ----
\newcommand{\R}{\mathbb{R}}
\newcommand{\grad}{\nabla}

% -------------------------------------------------------
\title{\textbf{Devoir d'Analyse de Régression}\\[0.4em]
       \large Calcul de Gradients et de Hessiennes -- Problème 1}
\author{Cerveau KATCHON}
\date{}
% -------------------------------------------------------

\begin{document}

\maketitle
\thispagestyle{fancy}

\section*{Problème 1 : Dérivées matricielles}

Ce problème porte sur le calcul du gradient vectoriel $\grad_x f(x)$ et de la
hessienne $\grad_x^2 f(x)$ de diverses fonctions scalaires de $x \in \R^n$.

% ==============================================================
\subsection*{Question 1 -- Gradient d'une forme quadratique généralisée}
% ==============================================================

\textbf{Énoncé.}
Soit la fonction
\[
  f(x) = \frac{1}{2}\,x^{\top} A x + b^{\top} x + c,
  \qquad
  A \in \R^{n \times n} \text{ symétrique},\quad
  x,\, b \in \R^{n},\quad
  c \in \R.
\]
Calculer $\grad_x f(x)$.

\medskip
\textbf{Démonstration.}

En appliquant la linéarité de la dérivée et les règles de dérivation matricielle :
\begin{align}
  \grad_x f(x)
    &= \frac{1}{2}\,\frac{\partial\bigl(x^{\top} A x\bigr)}{\partial x}
       + \frac{\partial\bigl(b^{\top} x\bigr)}{\partial x}
       + \frac{\partial c}{\partial x} \notag\\[6pt]
    &= \frac{1}{2}\,(A + A^{\top})\,x + b + 0 \notag\\[6pt]
    &= \frac{1}{2}\,2A\,x + b
       \qquad \bigl(\text{car } A^{\top} = A \text{ par hypothèse}\bigr) \notag\\[6pt]
    &= Ax + b. \notag
\end{align}

\medskip
\begin{center}
  \fbox{$\displaystyle \grad_x f(x) = Ax + b$}
\end{center}

% ==============================================================
\subsection*{Question 2 -- Hessienne d'une forme quadratique généralisée}
% ==============================================================

\textbf{Énoncé.}
Toujours avec $f(x) = \frac{1}{2}\,x^{\top} A x + b^{\top} x + c$
(A symétrique), calculer $\grad_x^2 f(x)$.

\medskip
\textbf{Démonstration.}

La hessienne est la dérivée du gradient par rapport à $x$ :
\begin{align}
  \grad_x^2 f(x)
    &= \frac{\partial^2 f(x)}{\partial x\,\partial x^{\top}}
     = \frac{\partial}{\partial x}\!\left(Ax + b\right)
     = A. \notag
\end{align}

\medskip
\begin{center}
  \fbox{$\displaystyle \grad_x^2 f(x) = A$}
\end{center}

% ==============================================================
\subsection*{Question 3 -- Gradient d'une exponentielle de forme quadratique}
% ==============================================================

\textbf{Énoncé.}
Soit $f(x) = e^{x^{\top} A x}$, $A \in \R^{n \times n}$ symétrique.
Calculer $\grad_x f(x)$.

\medskip
\textbf{Démonstration.}

Posons $g(x) = x^{\top} A x$.
La règle de dérivation en chaîne donne :
\begin{align}
  \grad_x f(x)
    &= e^{g(x)}\cdot \frac{\partial g(x)}{\partial x} \notag\\[6pt]
    &= e^{x^{\top} A x}\cdot (A + A^{\top})\,x \notag\\[6pt]
    &= e^{x^{\top} A x}\cdot 2Ax
       \qquad \bigl(\text{car } A^{\top} = A\bigr). \notag
\end{align}

\medskip
\begin{center}
  \fbox{$\displaystyle \grad_x f(x) = 2\,e^{x^{\top} A x}\,Ax$}
\end{center}

% ==============================================================
\subsection*{Question 4 -- Hessienne d'une exponentielle de forme quadratique}
% ==============================================================

\textbf{Énoncé.}
Toujours avec $f(x) = e^{x^{\top} A x}$, calculer $\grad_x^2 f(x)$.

\medskip
\textbf{Démonstration.}

En partant du gradient obtenu à la Question 3 et en appliquant la règle de
dérivation du produit :
\begin{align}
  \grad_x^2 f(x)
    &= 2\,\frac{\partial}{\partial x}\!\left(e^{x^{\top} A x}\,Ax\right)
    \notag\\[6pt]
    &= 2\!\left[
         \frac{\partial\bigl(e^{x^{\top} A x}\bigr)}{\partial x}\cdot (Ax)^{\top}
         + e^{x^{\top} A x}\cdot \frac{\partial(Ax)}{\partial x}
       \right] \notag\\[6pt]
    &= 2\!\left[
         2\,e^{x^{\top} A x}\,Ax\,(Ax)^{\top}
         + e^{x^{\top} A x}\,A
       \right] \notag\\[6pt]
    &= e^{x^{\top} A x}\!\left[
         4\,Ax\,x^{\top} A^{\top} + 2A
       \right]. \notag
\end{align}

Puisque $A$ est symétrique ($A^{\top} = A$) :

\medskip
\begin{center}
  \fbox{$\displaystyle
    \grad_x^2 f(x)
    = e^{x^{\top} A x}\!\left(4\,Ax\,x^{\top} A + 2A\right)
  $}
\end{center}

% ==============================================================
\subsection*{Question 5 -- Gradient de la norme euclidienne au carré de $Ax$}
% ==============================================================

\textbf{Énoncé.}
Soit $f(x) = \|Ax\|^2$, $A \in \R^{m \times n}$.
Calculer $\grad_x f(x)$.

\medskip
\textbf{Démonstration.}

On réécrit $f$ en développant la norme :
\[
  f(x) = \|Ax\|^2 = (Ax)^{\top}(Ax) = x^{\top} A^{\top} A\, x.
\]

Posons $B = A^{\top} A$ (matrice $n \times n$, symétrique positive). Alors
$f(x) = x^{\top} B x$ est une forme quadratique, et :
\begin{align}
  \grad_x f(x)
    &= \frac{\partial\bigl(x^{\top} B x\bigr)}{\partial x}
     = (B + B^{\top})\,x \notag\\[6pt]
    &= 2B\,x
       \qquad \bigl(\text{car } B^{\top} = (A^{\top}A)^{\top} = A^{\top}A = B\bigr)
    \notag\\[6pt]
    &= 2\,A^{\top} A\, x. \notag
\end{align}

\medskip
\begin{center}
  \fbox{$\displaystyle \grad_x f(x) = 2\,A^{\top} A\, x$}
\end{center}

% ==============================================================
\subsection*{Question 6 -- Hessienne de la norme euclidienne au carré de $Ax$}
% ==============================================================

\textbf{Énoncé.}
Soit $f(x) = \|Ax\|^2 = (Ax)^{\top}(Ax)$, $A \in \R^{m \times n}$.
Calculer $\grad_x^2 f(x)$.

\medskip
\textbf{Démonstration.}

On a établi à la Question~5 que $\grad_x f(x) = 2\,A^{\top} A\,x$.
La hessienne est la dérivée du gradient par rapport à $x$ :
\begin{align}
  \grad_x^2 f(x)
    &= \frac{\partial}{\partial x}\!\left(2\,A^{\top} A\, x\right)
     = 2\,A^{\top} A. \notag
\end{align}

\medskip
\begin{center}
  \fbox{$\displaystyle \grad_x^2 f(x) = 2\,A^{\top} A$}
\end{center}

% ==============================================================
\subsection*{Question 7 -- Gradient et Hessienne d'une fonction avec barrière logarithmique}
% ==============================================================

\textbf{Énoncé.}
Soit la fonction
\[
  f(x) = \frac{1}{2}\,x^{\top} A x - b^{\top} x
         - \mu_1 \log x_1 - \mu_2 \log x_2,
  \qquad
  x = \begin{pmatrix} x_1 \\ x_2 \end{pmatrix} \in \R^{2},\quad
  \mu_1,\, \mu_2 > 0.
\]

\textbf{(a)} Calculer $\grad_x f(x)$. \quad
\textbf{(b)} Calculer $\grad_x^2 f(x)$.

\medskip
\textbf{Démonstration -- partie (a).}

On décompose $f$ en trois termes :
\begin{itemize}
  \item $\frac{\partial}{\partial x}\!\left(\frac{1}{2}x^{\top} A x\right) = Ax$
        (avec $A$ symétrique, résultat de la Question~1).
  \item $\frac{\partial}{\partial x}\!\left(-b^{\top} x\right) = -b$.
  \item $\frac{\partial}{\partial x_i}\!\left(-\mu_i \log x_i\right)
        = -\dfrac{\mu_i}{x_i}$,
        donc $\dfrac{\partial}{\partial x}\!\left(-\mu_1\log x_1 - \mu_2\log x_2\right)
        = -\begin{pmatrix} \mu_1/x_1 \\ \mu_2/x_2 \end{pmatrix}$.
\end{itemize}

Par linéarité :
\[
  \grad_x f(x) = Ax - b - \begin{pmatrix} \mu_1/x_1 \\ \mu_2/x_2 \end{pmatrix}.
\]

\medskip
\begin{center}
  \fbox{$\displaystyle
    \grad_x f(x) = Ax - b - \begin{pmatrix} \mu_1/x_1 \\ \mu_2/x_2 \end{pmatrix}
  $}
\end{center}

\medskip
\textbf{Démonstration -- partie (b).}

La hessienne est la dérivée du gradient.
Les termes issus de $Ax - b$ donnent $A$.
Les termes logarithmiques contribuent :
\[
  \frac{\partial^2}{\partial x_i^2}\!\left(-\mu_i \log x_i\right)
  = \frac{\mu_i}{x_i^2},
  \qquad
  \frac{\partial^2}{\partial x_i\,\partial x_j}\!\left(-\mu_i \log x_i\right) = 0
  \text{ pour } i \neq j.
\]

Donc :
\[
  \grad_x^2 f(x)
  = A
    + \begin{pmatrix} \mu_1/x_1^2 & 0 \\ 0 & \mu_2/x_2^2 \end{pmatrix}
  = A + \operatorname{diag}\!\left(\frac{\mu_1}{x_1^2},\,\frac{\mu_2}{x_2^2}\right).
\]

\medskip
\begin{center}
  \fbox{$\displaystyle
    \grad_x^2 f(x)
    = A + \operatorname{diag}\!\!\left(\frac{\mu_1}{x_1^2},\,\frac{\mu_2}{x_2^2}\right)
  $}
\end{center}

% ==============================================================
\subsection*{Question 8 -- Direction de Newton}
% ==============================================================

\textbf{Énoncé.}
En partant de l'approximation de Taylor du second ordre de $f$ au voisinage
de $x_k$, déterminer la direction $p^*$ qui minimise cette approximation.
En déduire la mise à jour de Newton $x_{k+1}$.

\medskip
\textbf{Démonstration.}

Notons $g = \grad_x f(x_k)$ et $H = \grad_x^2 f(x_k)$.
L'approximation quadratique de $f(x_k + p)$ s'écrit :
\[
  f(x_k + p)
  \approx f(x_k) + g^{\top} p + \frac{1}{2}\,p^{\top} H\, p
  =: W(p).
\]

Pour minimiser $W(p)$ par rapport à $p$, on annule son gradient :
\begin{align}
  \frac{\partial W(p)}{\partial p}
    &= g + H\,p = 0 \notag\\[4pt]
  \Longrightarrow \quad
  H\,p^* &= -g. \notag
\end{align}

En supposant $H$ inversible ($H \succ 0$), on obtient la direction de Newton :
\[
  p^* = -H^{-1}\,g
      = -\bigl[\grad_x^2 f(x_k)\bigr]^{-1}\,\grad_x f(x_k).
\]

La mise à jour de Newton est donc :
\[
  x_{k+1} = x_k + p^* = x_k - \bigl[\grad_x^2 f(x_k)\bigr]^{-1}\,\grad_x f(x_k).
\]

\medskip
\begin{center}
  \fbox{$\displaystyle
    H(x_k)\,p^* = -\grad_x f(x_k)
    \qquad \Longrightarrow \qquad
    x_{k+1} = x_k - H(x_k)^{-1}\,\grad_x f(x_k)
  $}
\end{center}

\end{document}
